\section{Conclusion}
\label{sec:conclusion}

We evaluated a hybrid quantum temporal convolutional network for daily
cross-sectional equity ranking under a protocol built to make the quantum
contribution identifiable: six conditions trained through one code path, four rolling-origin folds with an embargo, five seeds, annualised risk
metrics with bootstrap intervals, Deflated Sharpe Ratios, backtest-overfitting
probabilities and family-wise error control. On \num{500} out-of-sample
trading days the trained circuit is indistinguishable from the same circuit
frozen at initialisation, from a parameter-matched classical surrogate, and
from the network with the circuit removed; no condition is distinguishable from
zero, and ten of the eleven have no break-even transaction cost at all. Two
one-sided tests place the trained-versus-frozen difference inside \num{0.03}
annualised Sharpe in both markets, and inside \num{0.10} when the circuit is
trained to the full budget without early stopping, which is equivalence rather than a failure
to reject; the contrast against the no-circuit control is resolved only to
\num{0.50}.

The circuit is not the weak point. It has no barren plateau to twelve qubits,
and its output ranking is robust to gate noise; a second ansatz family, with
re-uploading and an observable read from every wire, reproduces the ordering and bounds its
trained-versus-frozen difference at \num{0.25} Sharpe. What is
expensive is measurement: at the spread the trained models produce on a real
cross-section, preserving the ranking takes about \num{8192} shots per
security at the median cell, and more than \num{131072} at the narrowest,
rather than the \num{256} a scan on random inputs suggests. The weak point is the task, on which the achievable signal is
smaller than the variance introduced by the training draw. Earlier reports of
advantage on this data came from a 22-day window, single-seed model selection,
a split adjustment applied in the wrong direction in the public implementation,
and a classical encoder that already accounts for most of the reported gain.

Two artefacts are released with the paper (\url{https://github.com/ChiShengChen/qtcnn-trained-frozen-absent}): a batched formulation of the quantum
convolutional layer that reproduces the per-sample implementation to
double precision at a \numrange{33}{72}-fold speed-up, and the evaluation stack itself,
calibrated on synthetic data with known answers. Both are reusable on any
cross-sectional ranking task. We hope the second is used more often than the
first.
